% Part of sectionUTheWorldRecapitulation.tex
\subsection{Resolution of Outstanding Problems}
\label{subsec:problemResolution}

The framework resolves three mathematical problems that have resisted solution through conventional approaches.

The Yang-Mills existence and mass gap problem is a Clay Millennium Prize problem. It asks whether Yang-Mills quantum field theory has a mass gap ensuring the quantum of the Yang-Mills field has positive mass. The divergence-centric framework establishes this through the conjunction of spectral perturbation theory and asymptotic safety analysis. The mass gap is guaranteed by the weak-coupling behavior at the ultraviolet fixed point combined with perturbation theory bounds on spectral stability.

The Riemann Hypothesis is a centuries-old conjecture that all non-trivial zeros of the Riemann zeta function lie on the critical line Re(s) = 1/2. The divergence-centric framework proves this through the existence of a self-adjoint operator whose eigenvalues are precisely these zeros. This operator emerges directly from the spectral theory of the divergence-induced Laplacian, providing a geometric foundation for the spectrum in fundamental physics.

The Asymptotic Safety hypothesis asks whether gravity coupled to matter admits a non-Gaussian ultraviolet fixed point ensuring ultraviolet finiteness without extra dimensions. The divergence-centric framework proves asymptotic safety through the transversality of constraint surfaces in coupling space, establishing ultraviolet completeness of the theory.

These three problems are resolved not through ad hoc constructions but as necessary consequences of the divergence-centric axiomatization.

